\documentclass[acmsmall, screen]{acmart}

\settopmatter{printacmref=false}
\setcopyright{none}

\usepackage{eucal}

\usepackage{enumerate,xspace}
\usepackage[capitalise]{cleveref}

%% Bibliography style
\bibliographystyle{ACM-Reference-Format}
\citestyle{acmauthoryear}   %% For author/year citations

\begin{document}

\title[Constructing and Analyzing Pangenome Graphs]{Constructing and Analyzing Pangenome Graphs: \\ Annotated Bibliography}
\author{Anshuman Mohan}
\affiliation{
\department{Bowers CIS}
\institution{Cornell University}
\city{Ithaca}
\country{USA}
}
\email{amohan@cs.cornell.edu}


\maketitle

\section{Overview}

In the field of computational biology, there is much excitement about the
\emph{pangenome}. The traditional approach in genomics has long been to
compare a given sample to a single reference genome, but this has a
severe drawback: it is inherently overfitted to the reference.
In contrast, the pangenome allows for a many-to-many comparison of samples,
thus allowing a much richer exploration of the similarities and differences
between organisms.

Once constructed, a pangenome can be queried to yield results that help
identify aberrations in the human genome that may be the sources of diseases and
genetic proclivities. Used at its full potential, this technology could
radically change medicial diagnoses and treatment. This has a
direct tie to Sustainable Development Goal 3, ``Good Health and Well-Being''
\cite{sdg}.

Traditional genomic sequencing is famously expensive in terms of time and money,
and the pangenome graph is even more so. We do not wish to create a health
solution that is open exclusively to the rich and powerful,
so the \emph{efficient} construction of
pangenomes is an area of active development.
This document surveys a few existing methods and summarizes their approaches.
A more direct comparison, along with opportunities for a combination of these
approaches, will follow in future work.

\section{GFA Construction}

\subsection{Two Techniques}

The pangenome is typically represented as a graph, and the
Graphical Fragment Assembly (GFA) format \cite{gfa} has
emerged as the research standard.
Broadly, there are two techniques for the construction of these graphs.

The first is multiple sequence alignment (MSA).
There exists a dynamic programming algorithm to align $N$ genetic sequences,
but the algorithm runs in $O(L^N)$ (where $L$ is the length of the
sequences and $N$ the number of sequences), so it is typically only used
for pairwise alignment ($N=2$). There are ways to approximate the true result of
aligning $N$ sequences, but these methods are lossy and failure-prone.
\citet{poa} provide an answer via \emph{partial order mutliple sequence
alignment}: the MSA graph they create after the first pairwise dynamic
alignment can itself be aligned, as a first-class citizen,
by pairwise dynamic programming.
When a new sequence is aligned against the MSA, it is guaranteed to be compared
against each sequence already in the MSA. This leads to a lossless graphical
representation that is truly extensible.
Their solution is also faster to compute
than the state of the art: linear instead of exponential.

\citet{genome_alignment} proceed by conducting a survey of four graphical
representations of the pangenome. They find shortcomings in the construction
strategies of many of these, and propose a six-step process framework,
which they call the ABC to G-enome alignment, that applies to all the
representations. The basic idea is to think about alignable sequences in a
principled and transitive way. Their work is more of a guiding light than a
fully-formed solutions---some of the six steps are not fully fleshed out---but
this survey is a helpful reminder of what makes pangenome graphs similar to
one other.

\subsection{The Bleeding Edge}

We now look at two projects that have taken a step beyond GFAs.

\citet{minigraph} observe that there is an additional kind of lossiness
that pangenomes have not cared about so far: after folding a sample into a
reference graph and querying the graph, it is hard to go back to the
sample's original coordinate. Their solution is to introduce the
\emph{reference GFA} format (rGFA) that extends GFAs with three tags.
This gives rGFAs an unchanging and unambigious coordinate system that does not
change as the pangenome is grown, extended, and modified.

One of the tools surveyed by \citet{genome_alignment} is called Cactus.
\citet{progcactus} extend Cactus with a progressive aligner.
They instantiate a ``guide tree'' that gives Cactus a principled way to
break down a massive alignment task into smaller, more tractable, tasks.
The same tree is then used to stitch the answers back together.
While Cactus ``scales quadratically with the number of bases in the alignment
problem'' \cite{progcactus}, Progressive Cactus scales linearly.

\section{Variation graphs and ODGI}

The ODGI project takes an opinionated step away from GFAs,
but still maintains compatability. This project is our link to this
area, so let us study it for a moment.

\citet{variationgraphs} introduce \emph{variation graphs},
which are sequence graphs (themselves standard in pangenome studies) combined
with a set of paths through that graph that are possible in that
organism's population. Variation graphs are very general, not requiring
a directionality or an initial reference to build on.
For instance, they can represent duplication in the genome
using cycles in the graph.
The authors provide a technique for mapping a genome on to a variation graph
efficiently and accurately.

\citet{odgi} introduce ODGI
(Optimized Dynamic Genome/Graph Implementation),
which is (1) a way to represent pangenomes as variation graphs,
and (2) a suite of commands that analyze these graphs efficiently.
One of the contributions of the ODGI representation is that it does
significant preprocessing on the graph and maintains a rich amount of
information to hand. The set of nodes and the set of paths
are maintained at the top level, and further, every node maintains lists
of in- and out-edges, paths that traverse that node, and so on.
This unlocks the potential for various kinds of parallelization
(nodewise, pathwise, or edgewise), depending on the command we wish to optimize.
Note that ODGI's variation-graph-based representation is nonstandard;
the ODGI project offers conversion methods to and from the GFA format.


\section{Graph Analysis}

The works above represent general solutions at a low level, but
there also exist toolkits that try to serve as a ``complete package''.

A well-loved solution in this space is PPanGGolin, due to \citet{pangolin}.
It is used to study bacterial pangenomes. After constructing a pangenome
graph, PPanGGolin helps partition the nodes into sensible regions of similarity
and difference using an expectation-minimization algorithm.
By the end, a nodes of the new graph corresponds to a \emph{gene family},
which an edge between two nodes means that the two gene families
share at least one adjacent pair of genes.

\section{Pangenome Without a Graph}

A very interesting contribution comes from \citet{panx}, who present panX.
PanX does not actually build a pangenome graph, preferring to do a
protein alignment search using another tool called \textsc{diamond}.
The chief contribution is a web-based GUI that focusses on visualizing
pangenomes in a manner that is appealing and helpful for science.

\section{Efforts in other Domains}

As we have seen, there are two major steps to the pangenome puzzle:
construction and analysis. The conventional wisdom is that the two are
deeply connected because decisions made at one one level impact the other.

Can we take a different approach?
Even while staying ignorant of the underlying construction algorithm,
can we use an ML framework to analyze a pangenome graph?
It would be interesting to study the applicability of frameworks such as
DRNets and CLR-DRNets \cite{drnets, clrdrnets},
which have a track record of success in
solving complex tasks in unsupervised or weakly-supervised settings.
The frameworks have been deployed on
sudoku demixing and crystal structure phase mapping, and it would be
interesting to see if the analysis of a pangenome can also be
rephrased as a demixing problem.
DRNets and CLR-DRNets require the availability of a few logical rules
about the domain they are working in. As with the examples above, the human
pangenome has a number of well-studied maxims that hold true as
a form of prior knowledge. Perhaps DRNets or CLR-DRNets can exploit these to
construct rich graphs in the style of PPanGGolin, where a node
represents not just a genetic sequence but a gene family?

\bibliography{bib}

\end{document}